\scanpage{94}
\uline{équivaut} à celle de $\mathrm{cl}\,H^{2n}_!(X,F)-\mathrm{cl}\,H^{2n-1}_!(X,F)$ mod la $=$ chose. On peut supposer $X$ [affine et] lisse sur $\mathbb F_p$ \inserted{\deleted{\unsure{puremt}}\unsure{purement|variété} de dim $n$}, alors $H^{2n}_!(X,F)$ est dual de $H^0(X,\mathscr{H}om(F,\mathbb Q_\ell(n))=H^0(X,\check F)(n)$, \inserted{donc} purement de poids $2n$. D'autre part, $H^{2n-1}_!(X,F)$ est dual de $H^1(X,\mathscr{H}om(F,\mathbb Q_\ell(n)))$, et admet en général des poids $2n-1$ \uline{et} du poids $2n-2$.
\margin{\uline{peut se raffiner}, \uline{en regardant une} \uline{autre filtration} \uline{que par les} poids, \uline{savoir le niveau} \textemdash}
% slanted

\deleted{Une analyse} Si $F=\deleted{\unsure{k}}\mathbb Q_\ell(0)$, \deleted{alors} donc \deleted{\unsure{le cas}} si on s'occupe des fonctions $\zeta$ habituelles, la connaissance ``birationnelle grossière'' de $\zeta$ est \unsure{essentiellt|essentielt|essentiellement} \unsure{équivalent|équivalente} : celle de $\pi_0(\overline X)$ comme \unsure{frob-ensemble|Frob-ensemble}, et celle de Albanese de $X$ sur $\mathbb F_p$ \ldots\ldots
\margin{Modulo conjectures de Tate}
% slanted

Une analyse plus serrée \deleted{\unsure{des zéros|des} \unsure{et fon}} (plus fine que par le seul \uline{poids}, mais en utilisant aussi le niveau) permet de récupérer plus par la connaissance de $\zeta(X\deleted{\unsure{,F}})$ \inserted{(disons, en prenant $F=\mathbb Q_\ell$)} modulo les $\zeta(Y\deleted{\unsure{,G|,F}})$ ($\dim Y<n$) \deleted{poids\ $\unsure{\leq p}$}
% archive sheet number: 93
